%! Author = sriadityavedantam
%! Date = 3/25/26

\section{Perfect Sets and Connected Sets}

\lesson{26}{Friday 10 October 2025 9:10}{Day 26}

\begin{theorem}[Nested Compact Sets Property]{
    If $C_1\supset C_2\supset\cdots$ are all compact, non-empty, then
    \[
        \bigcap_{i=1}^\infty C_i\neq\varnothing.
    \]
}
    Suppose
    \[
        \bigcap_{i=1}^\infty C_i = \varnothing.
    \]
    Then
    \[
        C_1\cap\left(\bigcap_{i=2}^\infty C_i\right)=\varnothing.
    \]
    Thus
    \[
        C_1\subset\left(\bigcap_{i=2}^\infty C_i\right)^c=\bigcup_{i=2}^\infty C_i^c.
    \]
    Hence $\{C_i^c\}_{i=2}^\infty$ is an open cover of $C_1$.
    Since $C_1$ is compact, there is $N\in\n$ such that
    \[
        C_1\subset\bigcup_{i=2}^N C_i^c.
    \]
    Therefore
    \[
        C_1\cap\bigcap_{i=2}^N C_i=\varnothing.
    \]
    But since $C_1\supset C_2\supset\cdots\supset C_N$, we have
    \[
        \bigcap_{i=2}^N C_i = C_N.
    \]
    So
    \[
        C_N=\varnothing,
    \]
    a contradiction.
\end{theorem}

\begin{definition}[Cantor Set]
    Take
    \begin{gather*}
        C_1 \coloneqq [0,1],\\
        C_2 \coloneqq \left[0, \dfrac{1}{3}\right]\cup \left[\dfrac{2}{3}, 1\right],\\
        C_3 \coloneqq \left[0,\dfrac{1}{9}\right]\cup\left[\dfrac{2}{9},\dfrac{1}{3}\right]\cup\left[\dfrac{2}{3},\dfrac{7}{9}\right]\cup\left[\dfrac{8}{9},1\right].\\
    \end{gather*}
    Hence,
    \[
        C_{n+1} \coloneqq \dfrac{1}{3}C_n\cup\left(\dfrac{2}{3} + \dfrac{1}{3}C_n\right).
    \]
    All $C_i$ are compact since they are closed and bounded.
    Also $C_i\neq\varnothing$.
    Then
    \[
        \C \coloneqq \bigcap_{i=1}^\infty C_i\neq\varnothing.
    \]
\end{definition}

\begin{remark}
    \begin{enumerate}
        \item $\C = \dfrac{1}{3}\C\cup\left(\dfrac{2}{3} + \dfrac{1}{3}\C\right)$.
        \item Length of $[0,1]\setminus \C$ is
        \[
            \dfrac{1}{3} + \dfrac{2}{9} + \dfrac{4}{27} + \ldots = \sum_{n=0}^\infty \dfrac{1}{3}\left(\dfrac{2}{3}\right)^n = 1.
        \]
    \end{enumerate}
\end{remark}

\begin{definition}[Perfect]
    $S\subset\r$ is perfect if $S$ is closed and every point of $S$ is a limit point of $S$.
\end{definition}

\begin{remark}
    Every point of $S$ is a limit point of $S$ if and only if $S$ contains no isolated points.
\end{remark}

\begin{proposition}{
    The Cantor Set is perfect.
}
\end{proposition}

\begin{example}[Non-example and Example of Perfect Sets]
    \begin{enumerate}
        \item $\left\{\dfrac{1}{n} : n\in\n\right\}\cup\{0\}$ is closed, but not perfect.
        \item $[a,b]$ is perfect.
        \item The Cantor Set is perfect.
    \end{enumerate}
\end{example}